\documentclass[12pt,a4paper]{article}

% ============================================================
% PACKAGES
% ============================================================
\usepackage[utf8]{inputenc}
\usepackage[T1]{fontenc}
\usepackage{lmodern}
\usepackage[french]{babel}
\usepackage{geometry}
\usepackage{graphicx}
\graphicspath{{image/}}
\usepackage[table]{xcolor}
\usepackage{listings}
\usepackage{booktabs}
\usepackage{array}
\usepackage{fancyhdr}
\usepackage{titlesec}
\usepackage{enumitem}
\usepackage{mdframed}
\usepackage{tikz}
\usepackage{float}
\usepackage{microtype}
\usepackage{setspace}
\usepackage{hyperref}

% ============================================================
% MISE EN PAGE
% ============================================================
\geometry{
  headheight=14pt,
  top=2.6cm,
  bottom=2.8cm,
  left=2.8cm,
  right=2.8cm
}
\setstretch{1.12}

% ============================================================
% COULEURS
% ============================================================
\definecolor{primary}{RGB}{28, 90, 168}
\definecolor{primarylight}{RGB}{88, 145, 220}
\definecolor{secondary}{RGB}{232, 240, 252}
\definecolor{codebg}{RGB}{248, 250, 255}
\definecolor{codeframe}{RGB}{170, 195, 230}
\definecolor{success}{RGB}{34, 130, 90}
\definecolor{warning}{RGB}{200, 120, 40}
\definecolor{darkgray}{RGB}{45, 55, 72}
\definecolor{lightgray}{RGB}{95, 110, 140}

% ============================================================
% EN-TETE ET PIED DE PAGE
% ============================================================
\pagestyle{fancy}
\fancyhf{}
\fancyhead[L]{\small\sffamily\textcolor{primary!85!black}{Projet M1 GIL --- Gestion de flotte}}
\fancyhead[R]{\small\sffamily\textcolor{primary!85!black}{Univ. Rouen Normandie \textperiodcentered{} 2025--2026}}
\fancyfoot[C]{\small\sffamily\textcolor{primary!70!black}{\thepage}}
\renewcommand{\headrulewidth}{0.4pt}
\renewcommand{\headrule}{\hbox to\headwidth{\color{primary!35}\leaders\hrule height \headrulewidth\hfill}}
\renewcommand{\footrulewidth}{0pt}

% ============================================================
% STYLE DES TITRES
% ============================================================
\titleformat{\section}
  {\Large\sffamily\bfseries\color{primary}}
  {\thesection.}{0.55em}{}

\titleformat{\subsection}
  {\large\sffamily\bfseries\color{primary!85!black}}
  {\thesubsection.}{0.45em}{}

\titleformat{\subsubsection}
  {\normalsize\sffamily\bfseries\color{primary!75!black}}
  {\thesubsubsection.}{0.4em}{}
\titlespacing*{\section}{0pt}{1.6ex plus .2ex}{1ex plus .1ex}
\titlespacing*{\subsection}{0pt}{1.3ex plus .15ex}{0.7ex plus .1ex}

% ============================================================
% STYLE DU CODE
% ============================================================
\lstset{
  backgroundcolor=\color{codebg},
  frame=single,
  rulecolor=\color{codeframe},
  basicstyle=\ttfamily\footnotesize,
  keywordstyle=\color{primary}\bfseries,
  commentstyle=\color{lightgray}\itshape,
  stringstyle=\color{success},
  breaklines=true,
  showstringspaces=false,
  numbers=left,
  numberstyle=\tiny\sffamily\textcolor{lightgray},
  numbersep=6pt,
  tabsize=2,
  captionpos=b,
  aboveskip=\smallskipamount,
  belowskip=\smallskipamount,
  extendedchars=true,
  keepspaces=true,
  literate=
    {á}{{\'a}}1 {é}{{\'e}}1 {è}{{\`e}}1 {à}{{\`a}}1 {ù}{{\`u}}1 {ç}{{\c{c}}}1
    {ô}{{\^o}}1 {î}{{\^i}}1 {ê}{{\^e}}1 {û}{{\^u}}1 {â}{{\^a}}1
    {É}{{\'E}}1 {À}{{\`A}}1 {Ç}{{\c{C}}}1 {œ}{{\oe}}1 {Œ}{{\OE}}1
}

% ============================================================
% BOITE INFO
% ============================================================
\newmdenv[
  backgroundcolor=secondary,
  linecolor=primary,
  linewidth=1.1pt,
  leftline=true,
  rightline=false,
  topline=false,
  bottomline=false,
  innerleftmargin=11pt,
  innerrightmargin=11pt,
  innertopmargin=9pt,
  innerbottommargin=9pt,
  skipabove=0.6em,
  skipbelow=0.6em
]{infobox}

\hypersetup{
  colorlinks=true,
  linkcolor=primary!90!black,
  urlcolor=primarylight,
  citecolor=primary!85!black,
  pdfborderstyle={/S/U/W 0}
}

% ============================================================
% DOCUMENT
% ============================================================
\begin{document}

% ------------------------------------------------------------
% PAGE DE TITRE
% ------------------------------------------------------------
\begin{titlepage}
  \centering
  \vspace*{1cm}

  {\Large\textbf{Universit\'{e} de Rouen Normandie}\\[0.3cm]
  \large Master 1 -- G\'{e}nie Informatique et Logiciel}

  \vspace{1.5cm}
  {\color{primary!55}\rule{\linewidth}{0.5mm}}\\[0.45cm]

  {\Huge\sffamily\bfseries\color{primary} Projet Gestion de Flotte\\[0.35cm]}
  {\LARGE\sffamily\bfseries\color{primary!80!black} Rapport de Semaine 4}\\[0.35cm]
  {\large\sffamily\color{primary!65!black} Services Conducteurs \& Maintenance, Bruno, Datafaker}

  {\color{primary!55}\rule{\linewidth}{0.5mm}}\\[1.5cm]

  \begin{minipage}{0.45\textwidth}
    \begin{flushleft}
      \textbf{\'{E}tudiants :}\\
      Ahcene Amouchas \\
      Florian Pepin \\[0.5cm]
      \textbf{Encadrants :}\\
      Lydia \& Luc
    \end{flushleft}
  \end{minipage}
  \hfill
  \begin{minipage}{0.45\textwidth}
    \begin{flushright}
      \textbf{Ann\'{e}e universitaire :}\\
      2025 -- 2026 \\[0.5cm]
      \textbf{Semaine :}\\
      4 / 7
    \end{flushright}
  \end{minipage}

  \vfill
  {\large 4 avril 2026}
\end{titlepage}

\tableofcontents
\newpage

% ============================================================
\section{Introduction}
% ============================================================

Apr\`{e}s la mise en place de l'infrastructure (S2) et du microservice V\'{e}hicules (S3), la semaine~4 compl\`{e}te la plateforme avec les services \textbf{Conducteurs} et \textbf{Maintenance}, la communication \'{e}v\'{e}nementielle Kafka, la passerelle GraphQL et l'authentification Keycloak.

\subsection{Objectifs}

\begin{itemize}[leftmargin=*]
  \item \textbf{Service Conducteurs} --- CRUD, suivi des permis, assignations v\'{e}hicule--conducteur.
  \item \textbf{Service Maintenance} --- Planification, historique par v\'{e}hicule, alertes pr\'{e}ventives.
  \item \textbf{Kafka} --- Saga par chor\'{e}graphie pour les transitions d'\'{e}tat inter-services.
  \item \textbf{Gateway GraphQL} --- Agr\'{e}gation multi-services (Apollo Server).
  \item \textbf{Keycloak} --- JWT + r\^{o}les pour s\'{e}curiser REST et GraphQL.
  \item \textbf{Bruno} --- Collection de tests API (remplacement de Postman).
  \item \textbf{Datafaker} --- G\'{e}n\'{e}ration automatique de donn\'{e}es de test au d\'{e}marrage.
\end{itemize}

% ============================================================
\newpage
\section{Architecture de persistance}
% ============================================================

\begin{infobox}
\textbf{Approche :} Une instance PostgreSQL unique (Bitnami Helm chart) h\'{e}berge des bases s\'{e}par\'{e}es par service. L'isolation logique est assur\'{e}e au niveau base de donn\'{e}es, pas au niveau instance.
\end{infobox}

\subsection{Bases de donn\'{e}es}

Les quatre bases sont cr\'{e}\'{e}es automatiquement au d\'{e}marrage de PostgreSQL via un script \texttt{initdb} d\'{e}clar\'{e} dans les values Helm :

\begin{lstlisting}[language=SQL, caption=Script d'initialisation (values.yaml fleet-infra)]
CREATE DATABASE driver_db;
CREATE DATABASE maintenance_db;
CREATE DATABASE vehicle_db;
CREATE DATABASE keycloak_db;
\end{lstlisting}

\begin{itemize}
  \item \texttt{vehicle\_db}, \texttt{driver\_db}, \texttt{maintenance\_db} --- une base par microservice m\'{e}tier.
  \item \texttt{keycloak\_db} --- base d\'{e}di\'{e}e \`{a} Keycloak (au lieu de polluer la base \texttt{postgres} par d\'{e}faut).
\end{itemize}

Le sch\'{e}ma (DDL) est g\'{e}r\'{e} par Hibernate (\texttt{ddl-auto: update}) : pas de scripts SQL manuels \`{a} maintenir.

\subsection{Seeding automatique avec Datafaker}

Plut\^{o}t que des fichiers \texttt{data.sql} statiques, chaque service embarque un \textbf{DataSeeder} activ\'{e} par le profil Spring \texttt{dev}. Au d\'{e}marrage, si la table est vide, des donn\'{e}es r\'{e}alistes sont g\'{e}n\'{e}r\'{e}es via la biblioth\`{e}que \textbf{Datafaker} :

\begin{lstlisting}[language=Java, caption=Extrait du DataSeeder (vehicle-service)]
@Component
@Profile("dev")
public class DataSeeder implements CommandLineRunner {
    @Override
    public void run(String... args) {
        if (vehicleRepository.count() > 0) return;
        Faker faker = new Faker(Locale.FRANCE);
        for (int i = 0; i < 10; i++) {
            Vehicle v = new Vehicle();
            v.setPlateNumber(faker.regexify("[A-Z]{2}-[0-9]{3}-[A-Z]{2}"));
            v.setBrand(faker.options().option("Renault", "Peugeot", ...));
            // ...
            vehicleRepository.save(v);
        }
    }
}
\end{lstlisting}

\begin{table}[H]
\centering
\small
\begin{tabular}{|l|l|l|}
  \hline
  \rowcolor{primary!15}
  \textbf{Service} & \textbf{Donn\'{e}es g\'{e}n\'{e}r\'{e}es} & \textbf{D\'{e}tails} \\
  \hline
  vehicle-service & 10 v\'{e}hicules & Plaques, VIN, marques, kilom\'{e}trage al\'{e}atoires \\
  \hline
  driver-service & 8 conducteurs & Noms, emails, t\'{e}l\'{e}phones + 1 permis chacun \\
  \hline
  maintenance-service & 12 maintenances & Types, priorit\'{e}s, dates, co\^{u}ts pour les compl\'{e}t\'{e}es \\
  \hline
\end{tabular}
\caption{Donn\'{e}es g\'{e}n\'{e}r\'{e}es par Datafaker (profil \texttt{dev})}
\end{table}

Le profil \texttt{dev} est activ\'{e} via \texttt{SPRING\_PROFILES\_ACTIVE=dev} dans le template Helm (Kubernetes) et dans le \texttt{compose.yaml} (Docker). Les seeders sont \textbf{idempotents} : ils ne s'ex\'{e}cutent que si la table est vide.

% ============================================================
\newpage
\section{Microservice Driver}
% ============================================================

\begin{infobox}
\textbf{Objectif :} G\'{e}rer le cycle de vie des conducteurs : profils, permis de conduire, coordination avec le service V\'{e}hicules lors des assignations.
\end{infobox}

\subsection{Mod\`{e}le de donn\'{e}es}

Relation \textbf{1-N} entre \texttt{Driver} et \texttt{DriverLicense}. Le service impl\'{e}mente le \textit{Soft Delete} via \texttt{@SQLDelete} / \texttt{@SQLRestriction} Hibernate.

\begin{lstlisting}[language=Java, caption=Entit\'{e} Driver (extraits)]
@Entity @Table(name = "drivers")
@SQLDelete(sql = "UPDATE drivers SET deleted_at = NOW() WHERE id = ?")
@SQLRestriction("deleted_at IS NULL")
public class Driver {
    @Id @GeneratedValue(strategy = GenerationType.UUID)
    private UUID id;
    private String firstName, lastName, email, phone, employeeId;

    @Enumerated(EnumType.STRING)
    private DriverStatus status = DriverStatus.ACTIVE;

    @OneToMany(mappedBy = "driver", cascade = CascadeType.ALL)
    private List<DriverLicense> licenses = new ArrayList<>();
}
\end{lstlisting}

\subsection{Surveillance des permis (LicenseScheduler)}

Une t\^{a}che \texttt{@Scheduled} (cron quotidien) d\'{e}tecte les permis expirant sous 30 jours et publie des \'{e}v\'{e}nements Kafka (\texttt{LICENSE\_EXPIRING}, \texttt{LICENSE\_EXPIRED}).

\subsection{Endpoints REST}

\begin{table}[H]
\centering\small
\begin{tabular}{|l|l|p{7cm}|}
  \hline
  \rowcolor{primary!15}
  \textbf{M\'{e}thode} & \textbf{Endpoint} & \textbf{Description} \\
  \hline
  GET & \texttt{/api/drivers} & Liste (filtrable par statut) \\
  POST & \texttt{/api/drivers} & Cr\'{e}ation d'un profil \\
  GET & \texttt{/api/drivers/\{id\}} & D\'{e}tails par UUID \\
  PUT & \texttt{/api/drivers/\{id\}} & Mise \`{a} jour compl\`{e}te \\
  PATCH & \texttt{/api/drivers/\{id\}/status} & Changement de statut \\
  POST & \texttt{/api/drivers/\{id\}/licenses} & Ajout d'un permis \\
  GET & \texttt{/api/drivers/\{id\}/licenses} & Liste des permis \\
  DELETE & \texttt{/api/drivers/\{id\}} & Suppression logique \\
  \hline
\end{tabular}
\caption{Endpoints REST --- service Driver}
\end{table}

% ============================================================
\newpage
\section{Microservice Maintenance}
% ============================================================

\begin{infobox}
\textbf{Objectif :} Tra\c{c}abilit\'{e} compl\`{e}te des interventions techniques, de la planification \`{a} la cl\^{o}ture, avec d\'{e}tection automatique des retards.
\end{infobox}

\subsection{Mod\`{e}le de donn\'{e}es}

L'entit\'{e} \texttt{MaintenanceRecord} utilise trois \'{e}num\'{e}rations (\texttt{MaintenanceType}, \texttt{MaintenanceStatus}, \texttt{MaintenancePriority}). Le lien vers le v\'{e}hicule est un simple UUID --- pas de cl\'{e} \'{e}trang\`{e}re inter-bases (principe microservices).

\subsection{D\'{e}tection des retards}

Une t\^{a}che \texttt{@Scheduled} quotidienne (01h00) passe les maintenances en retard au statut \texttt{OVERDUE} et publie une alerte Kafka.

\subsection{Endpoints REST}

\begin{table}[H]
\centering\small
\begin{tabular}{|l|l|p{7cm}|}
  \hline
  \rowcolor{primary!15}
  \textbf{M\'{e}thode} & \textbf{Endpoint} & \textbf{Description} \\
  \hline
  GET & \texttt{/api/maintenance} & Liste (filtrable par v\'{e}hicule, statut, priorit\'{e}) \\
  POST & \texttt{/api/maintenance} & Planification d'une intervention \\
  GET & \texttt{/api/maintenance/vehicle/\{id\}} & Historique d'un v\'{e}hicule \\
  PUT & \texttt{/api/maintenance/\{id\}} & Mise \`{a} jour g\'{e}n\'{e}rale \\
  PATCH & \texttt{/api/maintenance/\{id\}/status} & Cl\^{o}ture (co\^{u}t, km, notes) \\
  \hline
\end{tabular}
\caption{Endpoints REST --- service Maintenance}
\end{table}

% ============================================================
\newpage
\section{Gateway GraphQL}
% ============================================================

\begin{infobox}
\textbf{Objectif :} Agr\'{e}ger les trois services m\'{e}tier derri\`{e}re une API GraphQL unifi\'{e}e (Apollo Server, Node.js).
\end{infobox}

\subsection{Architecture}

Le Gateway est d\'{e}ploy\'{e} sur le port 4000 et route les requ\^{e}tes via des clients REST (Axios) vers chaque service backend. Le JWT re\c{c}u du client est \textbf{forward\'{e}} automatiquement aux services.

\begin{itemize}
  \item \textbf{VehicleClient} $\rightarrow$ \texttt{http://vehicle-service:8080}
  \item \textbf{DriverClient} $\rightarrow$ \texttt{http://driver-service:8080}
  \item \textbf{MaintenanceClient} $\rightarrow$ \texttt{http://maintenance-service:8080}
\end{itemize}

Un m\'{e}canisme de retry (intercepteur Axios) g\`{e}re les erreurs r\'{e}seau transitoires au d\'{e}marrage (\texttt{ECONNREFUSED}, \texttt{ECONNRESET}).

\subsection{Agr\'{e}gation multi-services}

Les relations suivantes sont r\'{e}solubles en une seule requ\^{e}te GraphQL :

\begin{itemize}
  \item \textbf{V\'{e}hicule $\leftrightarrow$ Maintenance} --- historique technique depuis la fiche v\'{e}hicule.
  \item \textbf{V\'{e}hicule $\leftrightarrow$ Conducteur} --- assignation courante et profil du conducteur.
  \item \textbf{Maintenance $\leftrightarrow$ Conducteur} --- technicien assign\'{e} \`{a} une intervention.
\end{itemize}

% ============================================================
\newpage
\section{Architecture \'{e}v\'{e}nementielle (Saga)}
% ============================================================

\begin{infobox}
\textbf{Principe :} Saga par \textbf{chor\'{e}graphie} --- chaque service r\'{e}agit aux \'{e}v\'{e}nements des autres et publie ses propres r\'{e}sultats. Pas d'orchestrateur central.
\end{infobox}

\subsection{Format d'enveloppe (KafkaEventEnvelope)}

\begin{lstlisting}[language=java, caption=Enveloppe \'{e}v\'{e}nement unifi\'{e}e]
{
  "eventId": "625c87b9-...",
  "eventType": "MAINTENANCE_COMPLETED",
  "timestamp": 1775220800.293,
  "payload": { "vehicle_id": "39e75ab3-...", "cost_eur": 250.50 },
  "metadata": { "triggered_by": "maintenance-service" }
}
\end{lstlisting}

\subsection{Transactions compensatoires}

\begin{enumerate}
  \item \textbf{T1} --- Maintenance publie \texttt{MAINTENANCE\_STARTED}.
  \item \textbf{T2} --- V\'{e}hicule consomme et passe en \texttt{IN\_MAINTENANCE}.
  \item \textbf{C1} --- En cas d'\'{e}chec, V\'{e}hicule publie \texttt{MAINTENANCE\_REJECTED} $\rightarrow$ Maintenance annule la fiche.
\end{enumerate}

\subsection{Topics Kafka}

\begin{table}[H]
\centering\small
\begin{tabular}{|l|l|p{6cm}|}
  \hline
  \rowcolor{primary!15}
  \textbf{Topic} & \textbf{\'{E}v\'{e}nements} & \textbf{Effet} \\
  \hline
  \texttt{flotte.maintenance.events} & \texttt{STARTED}, \texttt{COMPLETED}, \texttt{CANCELLED} & V\'{e}hicule : \texttt{in\_maintenance} / \texttt{available} \\
  \hline
  \texttt{flotte.conducteurs.events} & \texttt{DRIVER\_CREATED/UPDATED}, \texttt{LICENSE\_EXPIRING/EXPIRED} & Notification inter-services \\
  \hline
  \texttt{flotte.assignments.events} & \texttt{VEHICLE\_ASSIGNED/UNASSIGNED} & Driver : \texttt{ON\_TOUR} / \texttt{ACTIVE} \\
  \hline
\end{tabular}
\caption{Topics Kafka et flux d'\'{e}v\'{e}nements}
\end{table}

% ============================================================
\newpage
\section{Authentification Keycloak}
% ============================================================

\begin{infobox}
\textbf{Approche :} Keycloak (v26.5.7) d\'{e}ploy\'{e} avec \texttt{--import-realm} pour charger automatiquement le realm \texttt{gestion-flotte}, ses clients et ses utilisateurs au d\'{e}marrage.
\end{infobox}

\subsection{Configuration}

\begin{itemize}
  \item \textbf{Base d\'{e}di\'{e}e} --- Keycloak utilise \texttt{keycloak\_db} au lieu de la base \texttt{postgres} par d\'{e}faut.
  \item \textbf{Realm} --- \texttt{gestion-flotte}, import\'{e} depuis \texttt{keycloak/realm-export.json}.
  \item \textbf{Client} --- \texttt{admin-cli} (public, \texttt{directAccessGrantsEnabled}).
  \item \textbf{Utilisateur} --- \texttt{test-user} (r\^{o}le \texttt{admin}), credentials hash\'{e}es (argon2) dans l'export.
\end{itemize}

\subsection{Flux d'authentification}

Les services Spring Boot valident les JWT via le \texttt{jwk-set-uri} de Keycloak. Le Gateway GraphQL forward le header \texttt{Authorization} re\c{c}u du client vers chaque service backend.

\begin{lstlisting}[language=bash, caption=R\'{e}cup\'{e}ration du token (grant type password)]
POST /auth/realms/gestion-flotte/protocol/openid-connect/token
  grant_type=password
  client_id=admin-cli
  username=test-user
  password=password
\end{lstlisting}

% ============================================================
\newpage
\section{Tests API avec Bruno}
\label{sec:tests-bruno}
% ============================================================

\begin{infobox}
\textbf{Migration Postman $\rightarrow$ Bruno :} la collection est d\'{e}sormais au format Bruno (fichiers \texttt{.bru} versionn\'{e}s dans Git), avec gestion automatique du token JWT via un script de collection.
\end{infobox}

\subsection{Organisation de la collection}

\begin{lstlisting}[language=bash, caption=Structure du dossier bruno/]
bruno/
  bruno.json                  # Métadonnées collection
  collection.bru              # Auth bearer + pre-request auto-token
  environments/
    Minikube.json             # Variables pour flotte.local (K8s)
    Docker.json               # Variables pour localhost (Compose)
  Auth/
    token.bru                 # Requête manuelle token Keycloak
  REST/
    service-vehicule/         # CRUD véhicules (5 requêtes)
    service-driver/           # CRUD conducteurs (8 requêtes)
    service-maintenance/      # CRUD maintenance (5 requêtes)
  GraphQL/
    service-vehicule/         # Query vehicles
    service-driver/           # Query drivers
    service-maintenance/      # Query maintenanceRecords
\end{lstlisting}

\subsection{Gestion automatique du token}

Le fichier \texttt{collection.bru} contient un script \texttt{pre-request} qui s'ex\'{e}cute avant chaque requ\^{e}te. Il r\'{e}cup\`{e}re automatiquement un JWT Keycloak et le met en cache jusqu'\`{a} expiration (5 minutes). Aucune action manuelle n'est requise.

\subsection{Environnements}

Deux environnements sont fournis, s\'{e}lectionnables dans Bruno :

\begin{itemize}
  \item \textbf{Minikube} --- \texttt{base\_url = http://flotte.local} (Ingress K8s).
  \item \textbf{Docker} --- \texttt{base\_url = http://localhost} (Docker Compose).
\end{itemize}

\subsection{Bilan des tests REST}

\begin{table}[H]
\centering\small
\begin{tabular}{|l|l|l|l|}
  \hline
  \rowcolor{primary!15}
  \textbf{Service} & \textbf{Endpoint} & \textbf{Code attendu} & \textbf{R\'{e}sultat} \\
  \hline
  Driver & POST /api/drivers & 201 & \textcolor{success}{\textbf{OK}} \\
  Driver & GET /api/drivers & 200 & \textcolor{success}{\textbf{OK}} \\
  Driver & GET /api/drivers/\{id\} & 200 & \textcolor{success}{\textbf{OK}} \\
  Driver & PUT /api/drivers/\{id\} & 200 & \textcolor{success}{\textbf{OK}} \\
  Driver & PATCH /api/drivers/\{id\}/status & 200 & \textcolor{success}{\textbf{OK}} \\
  Driver & POST /api/drivers/\{id\}/licenses & 201 & \textcolor{success}{\textbf{OK}} \\
  Driver & GET /api/drivers/\{id\}/licenses & 200 & \textcolor{success}{\textbf{OK}} \\
  Driver & DELETE /api/drivers/\{id\} & 204 & \textcolor{success}{\textbf{OK}} \\
  \hline
  Maintenance & POST /api/maintenance & 201 & \textcolor{success}{\textbf{OK}} \\
  Maintenance & GET /api/maintenance & 200 & \textcolor{success}{\textbf{OK}} \\
  Maintenance & GET /api/maintenance/vehicle/\{id\} & 200 & \textcolor{success}{\textbf{OK}} \\
  Maintenance & PUT /api/maintenance/\{id\} & 200 & \textcolor{success}{\textbf{OK}} \\
  Maintenance & PATCH /api/maintenance/\{id\}/status & 200 & \textcolor{success}{\textbf{OK}} \\
  \hline
  Vehicle & POST /api/vehicles & 201 & \textcolor{success}{\textbf{OK}} \\
  Vehicle & GET /api/vehicles & 200 & \textcolor{success}{\textbf{OK}} \\
  Vehicle & PUT /api/vehicles/\{id\} & 200 & \textcolor{success}{\textbf{OK}} \\
  Vehicle & POST /api/vehicles/\{id\}/assignments & 201 & \textcolor{success}{\textbf{OK}} \\
  Vehicle & GET /api/vehicles/\{id\}/assignments & 200 & \textcolor{success}{\textbf{OK}} \\
  \hline
\end{tabular}
\caption{Bilan des 18 endpoints REST test\'{e}s via Bruno}
\end{table}

\subsection{Bilan des tests GraphQL}

\begin{table}[H]
\centering\small
\begin{tabular}{|l|l|l|l|}
  \hline
  \rowcolor{primary!15}
  \textbf{Domaine} & \textbf{Op\'{e}ration} & \textbf{HTTP} & \textbf{R\'{e}sultat} \\
  \hline
  V\'{e}hicules & Query \texttt{vehicles} & 200 & \textcolor{success}{\textbf{OK}} \\
  Maintenance & Query \texttt{maintenanceRecords} & 200 & \textcolor{success}{\textbf{OK}} \\
  Conducteurs & Query \texttt{drivers} & 200 & \textcolor{success}{\textbf{OK}} \\
  \hline
\end{tabular}
\caption{Bilan des requ\^{e}tes GraphQL (Gateway, JWT via Bruno)}
\end{table}

% ============================================================
\newpage
\section{Tests automatis\'{e}s et couverture}
% ============================================================

\begin{infobox}
\textbf{Objectif :} Couverture minimale de \textbf{80\% en lignes} sur les tests unitaires (JaCoCo).
\end{infobox}

\subsection{Strat\'{e}gie}

\begin{itemize}
  \item \textbf{Tests unitaires} (JUnit 5, Mockito, \texttt{@WebMvcTest}) --- validation des services et endpoints.
  \item \textbf{Tests d'int\'{e}gration} (\texttt{@SpringBootTest}, H2) --- flux HTTP + persistance.
\end{itemize}

Le profil Maven \texttt{unit-coverage} isole les tests unitaires pour le rapport JaCoCo :

\begin{lstlisting}[language=bash, caption=G\'{e}n\'{e}ration du rapport JaCoCo]
./mvnw clean test -Punit-coverage
# Rapport : target/site/jacoco/index.html
\end{lstlisting}

\subsection{R\'{e}sultats JaCoCo}

\begin{figure}[H]
  \centering
  \includegraphics[width=\textwidth]{./image/driver-jacoco.png}
  \caption{JaCoCo --- \texttt{driver-service}}
  \label{fig:jacoco-driver}
\end{figure}

\begin{figure}[H]
  \centering
  \includegraphics[width=\textwidth]{./image/maintenance-jacoco.png}
  \caption{JaCoCo --- \texttt{maintenance-service}}
  \label{fig:jacoco-maintenance}
\end{figure}

\begin{figure}[H]
  \centering
  \includegraphics[width=\textwidth]{./image/vehicule-jacoco.png}
  \caption{JaCoCo --- \texttt{vehicle-service}}
  \label{fig:jacoco-vehicle}
\end{figure}

\begin{table}[H]
\centering\small
\begin{tabular}{|l|l|l|}
  \hline
  \rowcolor{primary!15}
  \textbf{Service} & \textbf{Couverture lignes} & \textbf{Objectif} \\
  \hline
  \texttt{driver-service} & $\approx$ 89\% & $\geq$ 80\% \\
  \texttt{maintenance-service} & $\approx$ 85\% & $\geq$ 80\% \\
  \texttt{vehicle-service} & $\approx$ 89\% & $\geq$ 80\% \\
  \hline
\end{tabular}
\caption{Synth\`{e}se couverture JaCoCo (profil \texttt{unit-coverage})}
\end{table}

\end{document}
